\documentclass[11pt,a4paper]{article}
\usepackage[utf8]{inputenc}
\usepackage[T1]{fontenc}
\usepackage[a4paper,margin=1in]{geometry}
\usepackage{amsmath,amssymb}
\usepackage{booktabs}
\usepackage{tabularx}
\usepackage{array}
\usepackage{longtable}
\usepackage{graphicx}
\usepackage{float}
\usepackage{caption}
\usepackage{xcolor}
\usepackage{enumitem}
\usepackage{fancyhdr}
\usepackage{tikz}
\usetikzlibrary{shapes.geometric,arrows.meta,positioning,fit,backgrounds,calc}
\usepackage[hidelinks]{hyperref}

\definecolor{accent}{HTML}{1F4E79}
\definecolor{soft}{HTML}{F2F4F7}
\definecolor{rule}{HTML}{B9C2CC}
\definecolor{ok}{HTML}{2E7D32}
\definecolor{bad}{HTML}{B3261E}

\pagestyle{fancy}
\fancyhf{}
\renewcommand{\headrulewidth}{0.4pt}
\fancyhead[L]{\small\itshape ThinkStack}
\fancyhead[R]{\small\itshape \leftmark}
\fancyfoot[C]{\small\thepage}

\setlist{itemsep=2pt, topsep=4pt}

\newcommand{\coverpage}[2]{%
\begin{titlepage}
\centering
{\large CHRIST (Deemed to be University), Bengaluru\par}
{\large Department of Computer Science\par}
\vspace{1.1cm}
{\Large\bfseries #1\par}
\vspace{0.35cm}
{\large Course: CSC481-5, Computer Science Project\par}
{\large Semester V\par}
\vspace{1.5cm}
{\Huge\bfseries ThinkStack\par}
\vspace{0.35cm}
{\large\itshape #2\par}
\vspace{1.5cm}
{\large\bfseries Submitted by\par}
\vspace{0.3cm}
\begin{tabular}{clc}
\toprule
\textbf{Sl. No.} & \textbf{Student Name} & \textbf{Register No.} \\
\midrule
1. & Aditya Mehta & 2440204 \\
2. & Jitvan Chadha & 2440222 \\
3. & Ritesh K R & 2440233 \\
\bottomrule
\end{tabular}
\vspace{1.2cm}

\begin{tabular}{ll}
\textbf{Guide:} & Dr. New Begin \\
\textbf{Class:} & 5BSC CS \\
\textbf{Repository:} & \texttt{github.com/get-thinkstack/ThinkStack} \\
\end{tabular}
\vfill
{\large July 2026\par}
\end{titlepage}
}

\newcommand{\evidence}[1]{%
\begin{center}
\fcolorbox{rule}{soft}{\parbox{0.93\linewidth}{\small #1}}
\end{center}}

\newcommand{\yes}{\textcolor{ok}{\textbf{Yes}}}
\newcommand{\no}{\textcolor{bad}{\textbf{No}}}

\begin{document}
\coverpage{Project Diary}{A Record of Decisions, Implementations and Corrections}

\newpage
\section*{About this Diary}
\addcontentsline{toc}{section}{About this Diary}
\markboth{About this Diary}{}

This diary records the development of ThinkStack from 5 June to 30 July 2026.
It is reconstructed from the repository itself: 133 commits, 42 merged pull
requests, 32 recorded architecture decisions, and the published releases. It is
not written from memory.

Two conventions are observed throughout. First, every quantity was measured or
counted, never estimated. Second, decisions that turned out to be wrong are
recorded alongside those that were right, because a diary that contains only
successes teaches nothing.

The work was carried out jointly by the three team members. Entries describe
what was decided and built rather than who performed each action.

\vspace{0.6cm}
\noindent\textbf{Team:} Aditya Mehta (2440204), Jitvan Chadha (2440222), Ritesh K R (2440233)

\newpage
\tableofcontents
\newpage

\section{Timeline at a Glance}
\markboth{Timeline}{}

\begin{center}
\small
\begin{tabularx}{\linewidth}{@{}llX@{}}
\toprule
\textbf{Date} & \textbf{Phase} & \textbf{Event} \\
\midrule
5 Jun & Sprint I & repository created \\
17 Jun & Sprint I & project renamed ScholarLens to ThinkStack; decision log started \\
19 Jun & Sprint I & desktop shell introduced \\
20 Jun & Sprint I & encryption implemented \\
23 Jun & Sprint I & LaTeX editor and auto-healing compiler \\
24 Jun & Sprint I & interface unified; Sprint I closes \\
7 Jul & Sprint II & download page \\
10 Jul & Sprint II & cross-platform builds and CI \\
21 Jul & Sprint II & hardware-aware model routing \\
22 Jul & Sprint II & in-application updates \\
25 Jul & Sprint II & test suite introduced; Rust startup diagnosis \\
27 Jul & Sprint II & first public release, v0.1.0 \\
29 Jul & Sprint II & v0.1.1, then v1.0.0 \\
30 Jul & Sprint II & v1.0.0 found not to start; root cause analysis and correction \\
30 Jul & Sprint II & TeX engine bundled; validated beta published \\
\bottomrule
\end{tabularx}
\end{center}

\section{Sprint I: Building the Application}

\subsection{5 to 16 June: foundations}

The repository was created on 5 June. The first two weeks established the
project structure, an initial working demonstration, and the first iteration of
the interface.

\subsection{17 June: naming, audit, and the decision log}

Three things happened on the same day that shaped everything afterwards.

The project was renamed from ScholarLens to ThinkStack. A backend audit
corrected a set of defects found by reading rather than by testing. Most
consequentially, an architecture decision log was started.

\evidence{\textbf{Decision recorded.} Run entirely on the user's device, with no
hosted language model. The reasoning: a researcher working with unpublished
results or institutionally restricted material cannot upload it, and a student
project has no budget for per-call inference. This single decision determined
local inference, quantised models, a frozen runtime, and eventually a bundled
typesetting engine.}

The log reached 32 entries by the end of the project. Its value was not
apparent at the time; it became apparent during the dependency audit in the
final week, when the reasoning behind earlier choices could be read rather than
reconstructed.

\subsection{19 to 21 June: desktop shell and encryption}

The application was wrapped in a desktop shell. The shell starts the backend as
a supervised child process and stops it when the window closes, so no orphaned
process is left holding a port.

Encryption was implemented with Argon2id for key derivation and AES-256-GCM for
the ciphertext.

\evidence{\textbf{Decision recorded.} Use authenticated encryption rather than
an unauthenticated mode. Tampering with a stored paper is then detected at
decryption instead of producing plausible but corrupted output, which for
research material is the difference between an error and a silent falsification.}

\subsection{23 to 24 June: the paper writer}

The LaTeX editor was implemented, with a compiler designed to be forgiving. A
small language model produces markup that is frequently almost correct, so the
compiler inserts missing package declarations, wraps bare fragments into
complete documents, and continues past recoverable errors. When a pass produces
no PDF, the failing environment is progressively neutralised and compilation is
retried, so one malformed figure does not cost the author the paper.

Sprint I closed on 24 June with every feature working from a source checkout.

\section{Sprint II: Making It Usable by Others}

\subsection{7 to 10 July: distribution}

A download page was produced, and cross-platform builds were automated. The
build configuration was made data-driven, so that adding an operating system is
an edit to a configuration file rather than to build logic.

\subsection{20 to 22 July: hardware awareness}

Model selection was made dependent on the machine rather than fixed.

\evidence{\textbf{Decision recorded.} Size the model against \emph{available}
memory rather than installed memory. A user has a browser and an editor open.
Sizing against total memory produces a process that is terminated by the
operating system partway through an analysis, which is worse than producing a
slightly weaker answer.}

A related decision concerned graphics hardware. The diagnosis reports the GPU
but enables offload only for one specific class of device.

\evidence{\textbf{Decision recorded.} Report the GPU, but treat all but CUDA
devices with sufficient memory as unusable. The inference library shipped in the
installer is a CPU build, so attempting offload elsewhere fails at model load.
Reporting a capability the runtime cannot use converts a working configuration
into a crash.}

\subsection{25 July: testing and the Rust rewrite}

An automated test suite was introduced. Startup diagnosis was moved from Python
into the desktop shell, written in Rust, because the Python implementation
imported a deep learning framework purely to determine whether a GPU existed,
which cost several seconds on every launch.

\evidence{\textbf{Decision recorded, with a limit.} Hardware diagnosis moved to
Rust because it was on the startup path and slow. Model discovery was measured
at approximately 9 milliseconds and was left in Python, because the same
argument does not apply to code that is already fast. The principle recorded was
to move work for a measured reason, not because a language seems faster.}

\subsection{27 to 29 July: first releases}

Version 0.1.0 was published on 27 July, the first release to build on all three
operating systems. Two defects had blocked it, both worth recording because
neither was reproducible locally.

\begin{itemize}[leftmargin=*]
  \item The Windows build failed at model download. A JSON tool on that platform
        emits carriage returns, which left an invalid character on the end of a
        URL. The transfer tool rejected it. Linux and macOS were unaffected, so
        the defect existed only on the platform none of the team used daily.
  \item The macOS build failed at code signing. An absent optional credential
        was passed as an empty string, so the signing tool attempted to import
        an empty certificate rather than skipping signing.
\end{itemize}

Version 1.0.0 followed on 29 July, alongside a decision to reduce the installer
size.

\evidence{\textbf{Decision recorded.} Ship only the smaller model and offer the
larger one with the user's consent, reusing weights the user may already have
through other local tools. This removed approximately one gigabyte from every
installer. Matching required a canonical naming key, because the same weights
are named differently by each runtime, and comparing filenames caused the
application to offer a download the user did not need.}

\subsection{30 July: the release that did not work}

\evidence{\textbf{Observed.} The published v1.0.0 installer, downloaded and run,
displayed a loading screen indefinitely. One tester waited approximately 200
seconds. Continuous integration had been green throughout.}

This is the most instructive day of the project.

\subsubsection{Diagnosis}

The startup path was instrumented before anything was changed, on the reasoning
that the failure could not be fixed while it was invisible. The instrumentation
identified the cause in a single run.

Three defects compounded:

\begin{enumerate}[leftmargin=*]
  \item The packaging framework installs the executable and its resources into
        different directories. The lookup searched only beside the executable,
        found nothing, and fell back to the development path of launching a
        system Python interpreter.
  \item The Linux application image exports an environment variable that
        redirects a Python interpreter's library search into the mounted image.
        A frozen interpreter ignores this; a system interpreter does not, and it
        terminated before running any project code.
  \item The loading screen retried without limit and had no error state, so a
        backend that died in under a second looked exactly like one that was
        slow.
\end{enumerate}

\subsubsection{What the same investigation uncovered}

\begin{itemize}[leftmargin=*]
  \item The embedding weights had never been placed in the bundle. The first
        document ingested on an installed copy attempted a network download, in
        an application whose entire premise is that it needs no network. The
        code's own documentation stated the opposite.
  \item The model directory was derived from the current working directory,
        which the application image changes before launching. This was incorrect
        on every operating system for every installed copy; it worked only when
        launched from a developer checkout.
  \item The window terminated approximately one second after appearing, from
        heap corruption in a graphics path of the system webview. Every backend
        check passed while this happened, because the backend was healthy and
        only the window was gone.
\end{itemize}

\subsubsection{The correction that mattered most}

Rather than continue fixing defects individually, every external dependency was
enumerated by walking all imports and every subprocess call across 51 backend
modules.

\evidence{\textbf{Result.} Exactly one unbundled runtime dependency existed: the
TeX engine. Everything else was already inside the bundle. Bundling that engine
closed the entire category rather than one instance of it.}

The pattern the audit exposed was that the absent TeX engine, the absent
embedding weights and the incorrect model path were the same defect: an
assumption that a user's machine resembles a developer's.

\subsection{30 July: bundling the typesetting engine}

The engine was measured before it was adopted: a 58 MB binary and a package
cache of approximately 47 MB, against roughly 1.1 GB of headroom under the
hosting platform's asset limit.

The first attempt worked locally and failed on the Linux build server, which
produced the most transferable lesson of the project.

\evidence{\textbf{Cause.} The chosen engine version required a newer system C
library than the build platform provides, so it never executed at all, failing
in 15 milliseconds. The build step recorded a warning and continued, so the
installer was assembled with an engine that had no package cache. The same
binary would have reached users of long-term-support distributions, who would
have had a paper writer that silently could not compile.}

\evidence{\textbf{Correction.} The engine version was chosen by measuring the
minimum system library each release requires, rather than by taking the newest
release. The selected version runs on the build platform and on older user
systems, and reduced the bundled payload from 104 MB to 70 MB. Equally
important, the build step now prints the engine's actual error and stops the
build, because suppressing that error is what made the diagnosis expensive.}

\subsection{30 July: verified delivery}

A validated beta was published to all three operating systems, with every
automated check passing, including compilation of a PDF by the packaged
application.

\section{Decisions and Their Consequences}

\begin{center}
\small
\begin{tabularx}{\linewidth}{@{}XXX@{}}
\toprule
\textbf{Decision} & \textbf{Reason} & \textbf{Consequence accepted} \\
\midrule
local inference only & privacy, cost, offline operation & far weaker models; seconds per response \\
quantised small models & fit consumer memory & reduced quality on structured output \\
frozen Python runtime & the user installs nothing & installers approach 1 GB \\
file-based vector store & no database to install & keyword index rebuilt per query \\
bundle only the smaller model & installer size & analysis quality depends on a consented download \\
bundle a typesetting engine & the writer must work with no LaTeX & 70 MB added per installer \\
CPU-only inference build & one build behaves identically everywhere & GPU hardware unused \\
tags, not branches, produce releases & releases cannot happen by accident & a deliberate step is required \\
updates only when requested & an offline product must not contact a server unprompted & a user who never asks never updates \\
\bottomrule
\end{tabularx}
\end{center}

\section{What Was Learned}

\subsection{Testing source code is not testing a product}

309 tests passed while the shipped installer could not start. The tests examined
a source tree; users run an installer. The verification tiers added afterwards
exercise the packaged application on each supported operating system, on
machines that have no model and no typesetting distribution installed, so that a
pass is evidence about a user's machine.

\subsection{Suppressed errors are expensive}

Two separate diagnoses were slowed by output sent to a null device. In both
cases the fix was to print the error and stop, rather than warn and continue.
A build that continues after a step has failed produces an artefact nobody
intended to ship.

\subsection{Fix categories, not instances}

Three defects were the same defect. Enumerating every dependency once bounded
the problem and ended it, where fixing them one at a time as testers reported
them would not have converged.

\subsection{Documentation drifts silently}

Instructions told macOS users to use a mechanism the operating system had
removed, and a LaTeX error message offered a command for one Linux distribution
to every user on every platform. Both were written when they were true. Neither
was noticed until a person followed them and could not proceed.

\section{Present State}

\begin{center}
\small
\begin{tabular}{@{}lr@{}}
\toprule
\textbf{Measure} & \textbf{Value} \\
\midrule
Commits & 133 \\
Pull requests merged & 42 \\
Application code (Python) & 8\,258 lines, 73 modules \\
Test code & 2\,709 lines, 23 modules, 309 tests \\
Interface code & 3\,280 lines, 9 components \\
Desktop shell (Rust) & 745 lines \\
REST endpoints & 36 \\
Architecture decisions recorded & 32 \\
Releases published & 5 \\
Platforms built and validated automatically & 3 \\
\bottomrule
\end{tabular}
\end{center}

\subsection{Verification status by platform}

\begin{center}
\small
\begin{tabularx}{\linewidth}{@{}lX@{}}
\toprule
\textbf{Platform} & \textbf{Status} \\
\midrule
Linux & automatically validated and exercised by hand on physical hardware \\
Windows & automatically validated and confirmed working by a human tester, though not yet across a range of machines \\
macOS & automatically validated; windowed behaviour on physical hardware remains unconfirmed, because unsigned applications are obstructed on first launch \\
\bottomrule
\end{tabularx}
\end{center}

\section{Planned Work}

The following are planned rather than aspirational. Each follows from a
limitation identified during development, and several are already partly
prepared for in the existing structure of the system.

\subsection{Language models trained for this application}

The clearest quality limitation is the capability of a general-purpose 0.5
billion parameter model on structured output. Asked to plot a function it
produced a reference to an image file that does not exist, where a 1.5 billion
parameter model produced a correct plot. The intended response is models suited
to the task rather than simply larger ones, in three stages of increasing
ambition.

\paragraph{Stage one: fine-tuned task models.} Training pairs of natural
language instruction and compilable LaTeX are already collected passively during
normal use, and the routing table in the inference client already contains
entries for a dedicated writing model, so adopting one is a matter of supplying
weights. A second model fine-tuned for the structured summarisation and claim
extraction formats the application parses would address the cases where a
general model returns prose where JSON was requested.

\paragraph{Stage two: adapter distribution.} Rather than shipping whole models,
distribute small adapters applied over the bundled base model, so a capability
upgrade costs tens of megabytes instead of gigabytes. Selection would use the
hardware profile the application already computes at startup, so a machine
receives only variants it can run.

\paragraph{Stage three: models trained from scratch.} The longer term intention
is to train small language models specifically for academic literature work,
rather than adapting general-purpose ones. A model trained from the outset on
research text, structured extraction and typesetting markup can be far smaller
than a general model of equivalent usefulness for these tasks, which matters
directly for a product constrained to consumer hardware.

\paragraph{Federated learning.} A possible direction, subject to careful design,
is improving these models from real usage without collecting anyone's documents.
Training would occur on the user's own machine and only model updates, never
text, would leave the device, and only with explicit consent. This is recorded
as a possibility rather than a commitment: it must not be allowed to weaken the
privacy guarantee that is the reason the product exists, and that constraint
takes precedence over the capability.

\subsection{An interface design philosophy, and the refactor that follows it}

The interface was built feature by feature under time pressure. It is coherent
but was not designed as a system: styling has accumulated into a single large
stylesheet with two competing styling mechanisms in use.

The intention is deliberately ordered. A design philosophy will first be agreed
by the team through discussion, and the interface will then be refactored to it.
Refactoring before that agreement would simply produce a second set of
accumulated decisions. Points to be settled include:

\begin{itemize}[leftmargin=*]
  \item \textbf{Design tokens.} Colour, spacing and typography expressed as
        variables, so the interface compiles from a defined system rather than
        from decisions made one screen at a time.
  \item \textbf{One environment rather than four.} Browsing, searching,
        interpreting and writing are currently separate screens. The intended
        direction is a single workspace whose density and visual quiet change
        with the task.
  \item \textbf{A single cross-platform identity} rather than per-platform
        conventions, so the application looks deliberate on every operating
        system.
  \item \textbf{Latency as a design material.} Local inference takes seconds and
        that cannot be removed. The interface should be designed around honest
        waiting rather than concealing it.
  \item \textbf{Reading comfort.} Sessions are long and the content is dense, so
        typography, measure and contrast are functional requirements rather than
        decoration.
\end{itemize}

\subsection{LitGraph: consolidating retrieval, analysis and conversation}

At present the gap finder, search, and summarisation and chat are separate
features reached from separate screens, although they operate on the same corpus
and much of the same computation. The planned direction is to consolidate them
into a single capability, provisionally called LitGraph, in which the literature
is treated as a connected structure rather than as a list of documents.

\begin{itemize}[leftmargin=*]
  \item A user would ask a question, see the passages that answer it, see how
        the papers containing them relate to one another, and see where the
        literature is thin, as one continuous activity rather than four.
  \item Retrieval, summarisation and gap detection already share the embedding
        index and the same model runtime, so consolidation removes duplicated
        computation as well as duplicated interface.
  \item Relationships between papers become first-class, which is the
        foundation the previously planned citation visualisation requires.
\end{itemize}

\subsection{A robust paper editor}

The paper writer is the most used feature and has received the most correction,
but its editing surface remains a plain text area over LaTeX source. The planned
work is a genuinely capable editor:

\begin{itemize}[leftmargin=*]
  \item syntax awareness, bracket matching, and compiler diagnostics shown
        against the line that caused them;
  \item a visual editing mode in which document structure can be edited without
        the author reading markup, with the compiled PDF remaining the single
        source of truth for appearance;
  \item position synchronisation between the source and the compiled document,
        so that selecting a place in one moves to it in the other;
  \item continued strengthening of the natural language to LaTeX path, which
        currently rewrites a selected passage in place and is the feature the
        planned writing model would most improve.
\end{itemize}

\subsection{Further planned refactors}

\begin{itemize}[leftmargin=*]
  \item \textbf{Incremental keyword index}, removing the per-query rebuild that
        currently limits corpus size.
  \item \textbf{Stylesheet consolidation}, resolving the two styling mechanisms
        into one, carried out as part of the interface refactor above rather
        than separately.
  \item \textbf{Code signing and notarisation} for macOS and Windows. This is
        required before a production release: both operating systems currently
        obstruct the first launch of an unsigned application, which every user
        encounters and which reads as the software being broken.
  \item \textbf{Bundled model integrity checks}, so that a corrupted or partial
        payload is detected at startup rather than at first use.
\end{itemize}

\section{Closing Note}

The technically demanding parts of this project, namely local inference with
hardware-aware model selection, hybrid retrieval, an auto-healing typesetting
pipeline and authenticated encryption, were completed and work. What determined
whether anyone else could use the software was packaging, dependency
completeness, and verification of the delivered artefact.

The distinction between software that works and software that has been shown to
work on a machine other than its author's is the lesson this project taught most
clearly, and it was taught by a release that failed rather than by one that
succeeded.

\end{document}
